\chapter{子列与聚点}
\label{chap:subsequences-cluster-points}
\label{chap:subsequences-bw}

发散数列的全部项未必靠近同一点，但某些按原次序选出的项仍可能收敛。子列就是这种选择的精确形式。部分极限记录子列能够达到的极限，聚点则从每个尾部是否反复进入邻域来描述同一现象。

\section{子列及其指标结构}


\begin{definition}
设 \((a_n)\) 为数列。若正整数列
\[
 n_1<n_2<\cdots
\]
严格递增，则 \((a_{n_k})_{k\ge1}\) 称为 \((a_n)\) 的一个\term{子列}。映射 \(k\mapsto n_k\) 是 \(\Npos\) 到 \(\Npos\) 的严格递增映射。
\end{definition}

严格递增保证两点：子列保持原来的先后次序，并且指标趋于无穷。事实上由归纳有 \(n_k\ge k\)。只要求 \(n_k\) 互异而不递增会允许重排；只要求非递减则可能无限重复某一固定项，不能反映原列尾部。

\begin{example}
偶数项 \((a_{2k})\)、平方指标项 \((a_{k^2})\) 都是子列。序列
\((a_1,a_1,a_1,\ldots)\) 一般不是子列，因为对应指标不严格递增。
\end{example}

\begin{proposition}[子列的子列]
若 \((a_{n_k})\) 是 \((a_n)\) 的子列，而 \((a_{n_{k_j}})\) 是前者的子列，则它仍是 \((a_n)\) 的子列。
\end{proposition}

\begin{proof}
\(j\mapsto k_j\) 与 \(k\mapsto n_k\) 都严格递增，其复合 \(j\mapsto n_{k_j}\) 仍严格递增。
\end{proof}

\section{子列对收敛与发散的刻画}

\subsection{子列继承原列极限}


\begin{theorem}
若 \(a_n\to a\)，则每个子列 \(a_{n_k}\to a\)。若 \(a_n\to+\infty\) 或 \(-\infty\)，相应结论也成立。
\end{theorem}

\begin{proof}
给定 \(\varepsilon>0\)，取 \(N\) 使 \(n\ge N\) 时
\(\abs{a_n-a}<\varepsilon\)。由于 \(n_k\ge k\)，当 \(k\ge N\) 时 \(n_k\ge N\)，故
\(\abs{a_{n_k}-a}<\varepsilon\)。无穷极限只需把误差条件换成阈值条件，证明相同。
\end{proof}

\begin{corollary}
若原列有收敛子列，则原列不可能趋于与该子列极限不同的有限值，也不可能趋于 \(+\infty\) 或 \(-\infty\)。
\end{corollary}

\begin{insightbox}[一个收敛子列说明不了整列]
存在一个收敛子列不保证原列收敛。例如 \((-1)^n\) 的偶数子列收敛到 \(1\)，原列却发散。要从子列恢复原列，需要额外条件，例如原列为 Cauchy 列或单调列。
\end{insightbox}

\subsection{用子列证明原列不收敛}


\begin{theorem}[子列判别的反面]
若存在两条子列分别收敛到 \(a,b\)，且 \(a\ne b\)，则原数列发散。
\end{theorem}

\begin{proof}
若原列收敛到 \(L\)，两条子列都应收敛到 \(L\)。极限唯一性迫使 \(a=L=b\)，矛盾。
\end{proof}

\begin{example}
对 \(a_n=(-1)^n+1/n\)，偶数子列趋于 \(1\)，奇数子列趋于 \(-1\)，所以原列发散。项 \(1/n\) 的趋零扰动不消除主振荡。
\end{example}

\begin{proposition}
实数列 \((a_n)\) 不收敛到 \(a\)，当且仅当存在 \(\varepsilon_0>0\) 与一条子列 \((a_{n_k})\)，使
\[
 \abs{a_{n_k}-a}\ge\varepsilon_0\qquad(k\in\Npos).
\]
\end{proposition}

\begin{proof}
若不收敛到 \(a\)，严格否定给出 \(\varepsilon_0>0\)，使每个尾部都有坏项。先选 \(n_1\) 为某个坏指标；选好 \(n_k\) 后，对尾部起点 \(n_k+1\) 选坏指标 \(n_{k+1}\)，便得到严格递增指标列。反向若存在这样的子列，则任意尾部都含误差至少为 \(\varepsilon_0\) 的项，故不能收敛到 \(a\)。
\end{proof}

该命题把量词否定转化为可操作的反例子列，是\chapterref{chap:function-limits}证明函数极限序列刻画时的主要构造。

\begin{theorem}[进一步子列判据]
实数列 \((a_n)\) 收敛到 \(L\)，当且仅当它的每条子列都含有一条进一步子列收敛到 \(L\)。
\end{theorem}

\begin{proof}
若 \(a_n\to L\)，每条子列本身已经收敛到 \(L\)，因而任取其进一步子列即可。

反之，若原列不收敛到 \(L\)，由上一命题存在
\(\varepsilon_0>0\) 和一条子列 \((a_{n_k})\)，使
\[
 \abs{a_{n_k}-L}\ge\varepsilon_0\qquad(k\in\Npos).
\]
这条子列的任何进一步子列仍满足同一下界，不可能收敛到 \(L\)，与假设矛盾。
\end{proof}

这一判据不要求原列有界。它把“控制完整尾部”换成一个等价的抽取原则，在函数极限、紧致性和各种弱收敛问题中会反复出现。

\section{Bolzano--Weierstrass 定理与单调子列定理}

\subsection{二分区间法证明 Bolzano--Weierstrass 定理}


\begin{theorem}[Bolzano--Weierstrass 定理]
每个有界实数列都存在收敛子列。
\end{theorem}

\begin{proof}
设所有项都位于闭区间 \(I_0=[A,B]\)。若 \(A=B\)，数列恒为 \(A\)，结论直接成立。以下设 \(A<B\)。

把 \(I_0\) 二分成两个闭半区间。至少一个半区间含无穷多个数列项；否则两边都只含有限多个指标，其并也只含有限多个指标，与数列有无穷多个指标矛盾。选定一个含无穷多项的半区间为 \(I_1\)。递归进行，得到
\[
 I_k=[\alpha_k,\beta_k],\qquad
 I_{k+1}\subseteq I_k,\qquad
 \beta_k-\alpha_k=2^{-k}(B-A),
\]
且每个 \(I_k\) 含无穷多个数列项。

由 Archimedes 性质，区间长度趋于零；区间套定理给出唯一交点 \(x\in\bigcap_kI_k\)。递归选指标 \(n_k>n_{k-1}\)，使
\(a_{n_k}\in I_k\)。这是可能的，因为 \(I_k\) 含无穷多个指标，删除前 \(n_{k-1}\) 个指标后仍有可选项。由于 \(a_{n_k},x\in I_k\)，
\[
 \abs{a_{n_k}-x}\le\beta_k-\alpha_k
 =2^{-k}(B-A)\longrightarrow0.
\]
故所得子列收敛到 \(x\)。
\end{proof}

“区间含无穷多个数列项”按指标计数，而不是按不同数值计数。若常值列所有项都等于 \(c\)，任何含 \(c\) 的区间都含无穷多个项，尽管值集只有一个元素。

\begin{corollary}
每个有界实数列至少有一个部分极限。每个无限有界实数集至少有一个聚点。
\end{corollary}

第二个结论先从集合中选取互异点形成数列，再对该数列应用定理。互异性保证极限点的每个邻域含集合中不同于它的点。

\subsection{单调子列定理及另一种证明}


\begin{definition}
若指标 \(n\) 满足
\[
 a_n\ge a_m\qquad(m>n),
\]
则称 \(n\) 为数列的一个\term{峰指标}。峰指标处的项不小于其后全部项。
\end{definition}

\begin{theorem}[单调子列定理]
每个实数列都有一个单调子列。
\end{theorem}

\begin{proof}
分两种情形。

若峰指标有无穷多个，把它们按递增顺序记为
\(n_1<n_2<\cdots\)。由峰指标定义，
\[
 a_{n_1}\ge a_{n_2}\ge\cdots,
\]
得到单调递减子列。

若峰指标只有有限多个，取 \(N\) 大于所有峰指标。对每个 \(n\ge N\)，\(n\) 不是峰指标，所以存在 \(m>n\) 使 \(a_m>a_n\)。从任意 \(n_1\ge N\) 开始，递归选
\[
 n_{k+1}>n_k,\qquad a_{n_{k+1}}>a_{n_k},
\]
得到严格递增子列。
\end{proof}

若原列有界，所得单调子列仍有界，因而由单调有界收敛定理收敛。这给出 Bolzano--Weierstrass 定理的第二种证明。反过来，用 Bolzano--Weierstrass 定理只能直接得到收敛子列，而收敛子列未必单调；单调子列定理包含不同的组合结构。

\section{数列值集的聚点与部分极限}

\subsection{数列值集的聚点与数列的部分极限}


\begin{definition}
设 \(E\subseteq\R\)。若对每个 \(\varepsilon>0\)，删心邻域
\[
 (x-\varepsilon,x+\varepsilon)\setminus\{x\}
\]
都含 \(E\) 中的点，则称 \(x\) 为 \(E\) 的\term{聚点}。
\end{definition}

\begin{definition}
若实数列 \((a_n)\) 有子列 \(a_{n_k}\to L\in\R\)，则称 \(L\) 为该数列的一个\term{部分极限}或\term{子列极限}。
\end{definition}

聚点属于集合的概念，部分极限属于带指标的数列。常值列 \(a_n=c\) 的值集是单点集，没有聚点；但 \(c\) 是该数列的部分极限。若一个数值只出现有限次，则它要成为部分极限，必须有其他数列值从附近逼近它。

\begin{proposition}
实数 \(L\) 是 \((a_n)\) 的部分极限，当且仅当对每个 \(\varepsilon>0\) 和每个 \(N\in\Npos\)，都存在 \(n\ge N\) 使
\[
 \abs{a_n-L}<\varepsilon.
\]
\end{proposition}

\begin{proof}
若 \(a_{n_k}\to L\)，给定 \(\varepsilon,N\)，取 \(k\) 足够大使
\(n_k\ge N\) 且 \(\abs{a_{n_k}-L}<\varepsilon\)。

反之，先在 \(N=1\)、误差 \(1\) 下选 \(n_1\)。选好 \(n_k\) 后，在尾部起点 \(n_k+1\)、误差 \(1/(k+1)\) 下选 \(n_{k+1}\)。于是指标严格递增且
\(\abs{a_{n_k}-L}<1/k\)，故子列收敛到 \(L\)。
\end{proof}

\subsection{部分极限的子列刻画}


上一命题也可写成：每个以 \(L\) 为中心的开区间与每个尾部值集
\[
 A_N=\{a_n:n\ge N\}
\]
都有非空交。为避免提前使用闭包语言，令
\[
\operatorname{adh}(A_N)
:=\{x\in\R:\text{\(x\) 的每个开区间邻域都与 \(A_N\) 相交}\}.
\]
若以 \(\mathcal L(a)\) 表示全部部分极限组成的集合，则
\[
 \mathcal L(a)=\bigcap_{N=1}^{\infty}\operatorname{adh}(A_N).
\]
\chapterref{chap:continuous-functions-closed-interval}建立闭集和闭包以后，
\(\operatorname{adh}(A_N)\) 将写成通常的 \(\overline{A_N}\)。

\begin{theorem}
若 \((a_n)\) 有界，则其部分极限集合非空且有界。若部分极限只有一个 \(L\)，则原列收敛到 \(L\)。
\end{theorem}

\begin{proof}
非空性由 Bolzano--Weierstrass 定理。每个子列仍受原列的共同界控制，序关系取极限后说明部分极限也位于同一闭区间内。

若原列不收敛到 \(L\)，可抽取始终与 \(L\) 保持固定正距离的子列。它仍有界，故有收敛的进一步子列；该新部分极限不同于 \(L\)，矛盾。
\end{proof}

\begin{proposition}[部分极限对极限运算封闭]
设 \((a_n)\) 有界，且 \(L_j\in\mathcal L(a)\)、\(L_j\to L\)。则
\(L\in\mathcal L(a)\)。
\end{proposition}

\begin{proof}
递归选取严格递增指标 \(n_j\)。因为 \(L_j\) 是部分极限，可在第 \(j\) 步从指标大于 \(n_{j-1}\) 的项中选取 \(a_{n_j}\)，使
\[
 \abs{a_{n_j}-L_j}<\frac1j.
\]
同时 \(L_j\to L\)，故
\[
 \abs{a_{n_j}-L}
 \le\abs{a_{n_j}-L_j}+\abs{L_j-L}\longrightarrow0.
\]
于是 \(a_{n_j}\to L\)，所以 \(L\) 也是部分极限。
\end{proof}

后文的闭集语言会把这一命题简称为“\(\mathcal L(a)\) 是闭集”。这里的对角选择揭示了真正使用的事实：第 \(j\) 步既要逼近变化中的目标 \(L_j\)，又要保证指标超过此前全部选择。

无界数列可能没有有限部分极限，例如 \(a_n=n\)；也可能同时有有限部分极限和无穷方向，例如 \(a_{2n}=0\)、\(a_{2n-1}=n\)。

\section{有界性、收敛性和序列紧致性的区别}


有界性控制数列所在的整体区间，收敛性要求尾部集中到单点。紧致性原型定理介于二者之间：
\[
 \text{收敛}\Longrightarrow\text{有界}
 \Longrightarrow\text{存在收敛子列},
\]
第二个箭头在实数中成立，而两个逆箭头均不成立。

\begin{counterexample}
\((-1)^n\) 有界且有收敛子列，但原列不收敛。数列
\[
 0,1,0,1/2,0,1/3,\ldots
\]
中，由非零项组成的子列趋于 \(0\)，由零项组成的子列也趋于 \(0\)，而且原列确实收敛到 \(0\)。所以，数列由两种取值模式交替组成并不能单独说明它发散；要看这些模式产生的子列极限是否一致。
\end{counterexample}

\begin{proposition}[有界性的抽取刻画]
实数列 \((a_n)\) 有界，当且仅当它的每条子列都含有收敛的进一步子列。
\end{proposition}

\begin{proof}
若原列有界，每条子列仍有界，Bolzano--Weierstrass 定理给出收敛的进一步子列。

反之，若原列无界，可以递归选择严格递增指标 \(n_k\)，使
\[
 \abs{a_{n_k}}>k.
\]
这条子列的每条进一步子列仍无界，因而不可能收敛，与假设矛盾。
\end{proof}

因此在实数中，收敛性可以拆成两件事：有界性保证每次抽取还能继续抽到收敛列；进一步子列判据再要求所有这些局部极限只能汇聚到同一点。进入一般度量空间后，前一种性质将发展为相对序列紧致性，且不再能简单地用有界性代替。

\section*{习题}
\addcontentsline{toc}{section}{习题}

\subsection*{A组　定义与基本方法}
\begin{enumerate}[label=\textbf{\thechapter.\arabic*.}]
 \exerciseitem{11.1} 证明严格递增正整数列满足 \(n_k\ge k\)，并推出 \(n_k\to+\infty\)。
 \exerciseitem{11.2} 判断 \(a_{2k}\)、\(a_{k^2}\)、\(a_{2-k}\)、\(a_{\lfloor\sqrt k\rfloor+1}\) 哪些构成子列，并说明理由。
 \exerciseitem{11.3} 证明子列的子列仍是原列子列。
 \exerciseitem{11.4} 从定义证明有限极限和无穷极限都由原列传给子列。
 \exerciseitem{11.5} 证明：若原列的每个子列都有进一步子列收敛到 \(L\)，则原列收敛到 \(L\)。
 \exerciseitem{11.6} 把“不收敛到 \(L\)”的量词否定递归构造成始终远离 \(L\) 的子列。
 \exerciseitem{12.1} 证明部分极限的“每个邻域命中任意尾部”刻画。
 \exerciseitem{12.2} 分别求 \((-1)^n\)、\((-1)^n+1/n\)、\(\sin(n\pi/2)\) 的全部部分极限。
 \exerciseitem{12.3} 给出数列部分极限但不是其值集聚点的例子，并解释差异。
 \exerciseitem{12.4} 证明有界数列的部分极限集合非空且有最大元、最小元。
\end{enumerate}

\subsection*{B组　证明与反例}
\begin{enumerate}[label=\textbf{\thechapter.\arabic*.},resume]
 \exerciseitem{11.7} 用奇偶子列判断 \((-1)^n+1/n\) 的收敛性。
 \exerciseitem{11.8} 补全二分区间证明中“至少一个半区间含无穷多项”和“可以选 \(n_k>n_{k-1}\)”两步。
 \exerciseitem{11.9} 证明任意有界数列都有 Cauchy 子列，并给出第 \(k\) 步的距离估计。
 \exerciseitem{11.10} 从无限有界集合中抽取互异数列，证明该集合有聚点。
 \exerciseitem{11.11} 证明峰指标无穷时得到递减子列，峰指标有限时得到严格递增子列。
 \exerciseitem{11.12} 用单调子列定理重新证明 Bolzano--Weierstrass 定理。
 \exerciseitem{11.13} 证明：有界数列若只有一个部分极限，则原列收敛。
 \exerciseitem{11.14} 给出无界数列具有收敛子列的例子，并说明这不与 Bolzano--Weierstrass 定理矛盾。
 \exerciseitem{11.15} 证明若数列的值只取有限多个实数，则至少有一个常值子列；刻画原列何时收敛。
\end{enumerate}

\subsection*{C组　综合问题}
\begin{enumerate}[label=\textbf{\thechapter.\arabic*.},resume]
 \exerciseitem{11.16} 构造有界数列，使其部分极限恰为给定的两个实数 \(a<b\)。
 \exerciseitem{11.17} 构造有界数列，其部分极限集合为 \([0,1]\)。允许使用有理数在区间中的可数稠密枚举。
 \projectexerciseitem{11.18} \ 【完成\chapterref{chap:equivalent-completeness}后回做；预计用时：75分钟】完整重建 Bolzano--Weierstrass 定理的二分区间证明，并给出子列误差不超过 \(2^{-k}(B-A)\) 的定量版本。
 \optionalexerciseitem{11.19} 证明任意实数列都有单调子列，并讨论在偏序集值数列中该结论为何可能失败。
 \optionalexerciseitem{11.20} 若 \(a_{n+1}-a_n\to0\)，两条收敛子列的极限是否必须相同？给出反例。
 \optionalexerciseitem{11.21} 证明有界数列收敛到 \(L\) 当且仅当每条子列都含有进一步子列收敛到 \(L\)。
 \openexerciseitem{11.22} 比较“值集有聚点”与“数列有部分极限”。说明重复取值为何使两概念不能简单等同。
 \projectexerciseitem{11.23} \ 【前置：\chapterref{chap:real-number-system}可数性；预计用时：60分钟】给定可数闭集的若干具体例子，设计数列使其部分极限集合等于该集合；分析构造中每个点需出现的逼近频率。
 \exerciseitem{12.20} 构造数列，使其部分极限集合恰为 \(\{-1,0,2\}\)，并计算上下极限。
 \projectexerciseitem{12.22} \ 【前置：\chapterref{chap:subsequences-cluster-points}Bolzano--Weierstrass；预计用时：90分钟】证明有界数列的部分极限集合是非空闭集，且其最大、最小元分别为上下极限。
 \openexerciseitem{12.25} 研究任意非空闭集 \(F\subseteq\R\) 是否能成为某个数列的部分极限集合；对有界 \(F\) 给出构造方案。
\end{enumerate}

\section*{部分习题提示}
\begin{itemize}
 \item \exerciseref{11.5}：反设不收敛，使用\exerciseref{11.6}构造坏子列。
 \item \exerciseref{11.13}：反设不收敛到唯一部分极限，再抽取远离它的子列并应用 Bolzano--Weierstrass。
 \item \exerciseref{11.17}：把 \(\Q\cap[0,1]\) 的前 \(m\) 个数依次组成第 \(m\) 个有限块。
 \item \exerciseref{11.20}：可使用步长逐渐缩小、往返于两个端点的三角形块。
\end{itemize}
